\begin{trapbox}

	\begin{description}[style=nextline, leftmargin=2.8em, labelsep=0.5em, itemsep=0.35em]
		\item[\textbf{\textcolor{CTrap}{易错点一（焦点混淆）：}}]
			当双曲线焦点在 $y$ 轴时，错误地使用焦点在 $x$ 轴的公式 $\frac{x_0 x}{a^2} - \frac{y_0 y}{b^2}=1$。
		\item[\textbf{\textcolor{CTrap}{正确理解（区分焦点）：}}]
			焦点在 $y$ 轴时双曲线方程为 $\frac{y^2}{a^2} - \frac{x^2}{b^2}=1$，对应的切点弦方程为 $\frac{y_0 y}{a^2} - \frac{x_0 x}{b^2}=1$，注意 $y_0 y$ 项放在分子前面且分母为 $a^2$。
		\item[\textbf{\textcolor{CTrap}{错因分析（记忆混乱）：}}]
			未根据焦点位置调整公式，死记硬背形式，导致 $x$ 与 $y$ 位置颠倒。
	\end{description}

	\medskip
	\begin{description}[style=nextline, leftmargin=2.8em, labelsep=0.5em, itemsep=0.35em]
		\item[\textbf{\textcolor{CTrap}{易错点二（位置忽视）：}}]
			不对点 $P$ 是否在双曲线外部进行验证，直接使用切点弦公式。
		\item[\textbf{\textcolor{CTrap}{正确理解（使用前提）：}}]
			公式仅在点 $P$ 位于双曲线外部且可引两条实切线时成立。必须按上述 $q$ 条件检查；不能只凭含混的“内、外部”称呼作判断。
		\item[\textbf{\textcolor{CTrap}{错因分析（忽略条件）：}}]
			未注意定理前提，机械套用，导致结论错误。
	\end{description}

	\medskip
	\begin{description}[style=nextline, leftmargin=2.8em, labelsep=0.5em, itemsep=0.35em]
		\item[\textbf{\textcolor{CTrap}{易错点三（符号错误）：}}]
			切点弦方程中间符号写为加号 $+$，误用椭圆中的规律。
		\item[\textbf{\textcolor{CTrap}{正确理解（保持减号）：}}]
			双曲线方程中间是减号，切点弦方程也是减号，形式为 $\frac{x_0 x}{a^2} - \frac{y_0 y}{b^2}=1$ 或 $\frac{y_0 y}{a^2} - \frac{x_0 x}{b^2}=1$，不可改为加号。
		\item[\textbf{\textcolor{CTrap}{错因分析（类比混淆）：}}]
			椭圆切点弦方程为加号，学生类比时默认改成加号，未注意双曲线符号差异。
	\end{description}

\textbf{两有限实切点的存在条件。}
记 $q=x_0^2/a^2-y_0^2/b^2$，则恰有两个不同的有限实切点当且仅当
$q<1$ 且 $q\ne0$。$q=1$ 是一个切点，$q>1$ 无实切点；
$q=0$ 的非中心点在渐近线上，只有一个有限切点，另一方向为渐近线；中心点没有实切线。
纵轴形式使用 $q=y_0^2/a^2-x_0^2/b^2$。
\end{trapbox}
